

This section defines the maintenance and support requirements for the Agrihome smart plant monitoring system after deployment. The system consists of a camera module, processing unit, and backend software used to analyze plant images and provide health insights. Maintenance requirements ensure that the system remains reliable, updatable, and easy to troubleshoot over time. These requirements address software updates, hardware servicing, system monitoring, and documentation to support both developers and end users in maintaining proper system operation.

\subsection{System Maintenance Documentation}
\subsubsection{Description}
The system shall provide comprehensive maintenance documentation, including installation guides, system architecture descriptions, and troubleshooting procedures. This documentation shall enable developers and maintainers to understand system behavior, diagnose issues, and perform updates or repairs efficiently.

\subsubsection{Source}
Project design requirements and maintainability best practices.

\subsubsection{Constraints}
Documentation must be updated as system changes occur. Limited development time may restrict the level of detail in the initial release.

\subsubsection{Standards}
Documentation should follow standard software engineering documentation practices.

\subsubsection{Priority}
High


\subsection{Hardware Maintenance and Replacement}
\subsubsection{Description}
The system shall allow easy access to hardware components such as the camera module and processing unit for repair or replacement. Components should be modular and replaceable without requiring complete system disassembly.

\subsubsection{Source}
Hardware design considerations and user maintenance needs.

\subsubsection{Constraints}
Physical design and mounting structure may limit accessibility. Replacement components must be compatible with the system.

\subsubsection{Standards}
Hardware components should follow manufacturer specifications and safe handling practices.

\subsubsection{Priority}
Critical


\subsection{Software Updates and Bug Fixes}
\subsubsection{Description}
The system shall support software updates to fix bugs, improve performance, and enhance functionality. Updates should be deployable without requiring complete system reinstallation.

\subsubsection{Source}
Software lifecycle management requirements.

\subsubsection{Constraints}
Update process may depend on network availability. Improper updates may impact system stability.

\subsubsection{Standards}
Software updates should follow version control and deployment best practices.

\subsubsection{Priority}
Critical


\subsection{System Monitoring and Error Logging}
\subsubsection{Description}
The system shall maintain logs of system activity, errors, and failures to assist in troubleshooting. Logs should include timestamps and relevant system events to help diagnose issues efficiently.

\subsubsection{Source}
System reliability and debugging requirements.

\subsubsection{Constraints}
Storage limitations may restrict log retention. Logging must not significantly degrade system performance.

\subsubsection{Standards}
Logging should follow common software monitoring and debugging practices.

\subsubsection{Priority}
High


\subsection{Technical Support Availability}
\subsubsection{Description}
The system shall provide access to support resources such as user guides, FAQs, and developer documentation to assist users in resolving issues and understanding system functionality.

\subsubsection{Source}
Customer support and usability requirements.

\subsubsection{Constraints}
Support availability may be limited due to team size. Real-time support may not be feasible in the prototype phase.

\subsubsection{Standards}
Support materials should follow standard customer support documentation practices.

\subsubsection{Priority}
Medium